The disk and hole exchange their signs of growth between $q=0.0458221$
and 0.0460979. The primary mean disk rates differ from the integrated
curvature prediction by at most $2.65\times10^{-4}$ sites per step.
For holes the difference is at most $3.60\times10^{-4}$, except at
$q=0.0456842$, where it is $5.78\times10^{-4}$ and the surrounding
metastable active phase also begins to decay. Replacing the trimmed,
rectified area by the full signed area changes the mean hole rate at
$q=0.0458221$ from $0.710\times10^{-3}$ to $0.873\times10^{-3}$,
without changing its sign. The complementary growth signs persist in
every individual run at all four sampled probabilities.

These coarse droplet measurements locate opposing invasion tendencies.
They do not determine whether an asymptotic interval of generic
coexistence lies between the closest points. Even with a single isotropic
front reversal, the curvature law makes both droplets of a finite radius
shrink whenever $|v|<D/R$. Consequently, observing both finite droplets
shrink would also require a radius sequence before it could distinguish
such an interval from ordinary curvature. Here the central droplet gap
is $2.758\times10^{-4}$, well above the observed front-crossing differences.

